\ifdefined\pdfinfoomitdate
  \pdfinfoomitdate=1
\fi
\ifdefined\pdfsuppressptexinfo
  \pdfsuppressptexinfo=-1
\fi
\ifdefined\pdftrailerid
  \pdftrailerid{}
\fi

\usepackage[T1]{fontenc}
\usepackage[utf8]{inputenc}
\usepackage{textcomp}
\usepackage{fancyvrb}
\usepackage{framed}
\makeatletter
\@ifundefined{Shaded}{\newenvironment{Shaded}{}{}}{}
\makeatother
\ExplSyntaxOn
\cs_if_exist:NF \tbl_save_outer_table_cols:
  { \cs_new_protected:Npn \tbl_save_outer_table_cols: {} }
\ExplSyntaxOff
\usepackage{mathtools}
\usepackage{iftex}
\usepackage{enumitem}
\usepackage{longtable}
\usepackage{booktabs}
\usepackage{array}
\usepackage{colortbl}

\renewcommand{\bfdefault}{sb}

\renewcommand{\labelitemii}{\textopenbullet}
\renewcommand{\labelitemiii}{\raisebox{0.2ex}{\scriptsize$\blacksquare$}}
\renewcommand{\labelitemiv}{\raisebox{0.2ex}{\scriptsize$\blacksquare$}}

\setlist[enumerate,1]{label=\arabic*.,leftmargin=*}
\setlist[enumerate,2]{label=\alph*.,leftmargin=*}
\setlist[enumerate,3]{label=\roman*.,leftmargin=*}
\setlist[enumerate,4]{label=\roman*.,leftmargin=*}
\setlist[itemize,2]{label=\textopenbullet}
\setlist[itemize,3]{label=\raisebox{0.2ex}{\scriptsize$\blacksquare$}}
\setlist[itemize,4]{label=\raisebox{0.2ex}{\scriptsize$\blacksquare$}}

\ifPDFTeX
  % Arrows
  \DeclareUnicodeCharacter{2190}{\ensuremath{\leftarrow}}
  \DeclareUnicodeCharacter{2191}{\ensuremath{\uparrow}}
  \DeclareUnicodeCharacter{2192}{\ensuremath{\rightarrow}}
  \DeclareUnicodeCharacter{2193}{\ensuremath{\downarrow}}
  \DeclareUnicodeCharacter{2194}{\ensuremath{\leftrightarrow}}
  \DeclareUnicodeCharacter{2195}{\ensuremath{\updownarrow}}
  \DeclareUnicodeCharacter{2196}{\ensuremath{\nwarrow}}
  \DeclareUnicodeCharacter{2197}{\ensuremath{\nearrow}}
  \DeclareUnicodeCharacter{2198}{\ensuremath{\searrow}}
  \DeclareUnicodeCharacter{2199}{\ensuremath{\swarrow}}
  \DeclareUnicodeCharacter{21D0}{\ensuremath{\Leftarrow}}
  \DeclareUnicodeCharacter{21D2}{\ensuremath{\Rightarrow}}
  \DeclareUnicodeCharacter{21D4}{\ensuremath{\Leftrightarrow}}
  \DeclareUnicodeCharacter{21A6}{\ensuremath{\mapsto}}
  % Math operators and relations
  \DeclareUnicodeCharacter{00B1}{\ensuremath{\pm}}
  \DeclareUnicodeCharacter{00B7}{\ensuremath{\cdot}}
  \DeclareUnicodeCharacter{00D7}{\ensuremath{\times}}
  \DeclareUnicodeCharacter{00F7}{\ensuremath{\div}}
  \DeclareUnicodeCharacter{2200}{\ensuremath{\forall}}
  \DeclareUnicodeCharacter{2203}{\ensuremath{\exists}}
  \DeclareUnicodeCharacter{2205}{\ensuremath{\emptyset}}
  \DeclareUnicodeCharacter{2208}{\ensuremath{\in}}
  \DeclareUnicodeCharacter{2209}{\ensuremath{\notin}}
  \DeclareUnicodeCharacter{220B}{\ensuremath{\ni}}
  \DeclareUnicodeCharacter{2211}{\ensuremath{\sum}}
  \DeclareUnicodeCharacter{2212}{\ensuremath{-}}
  \DeclareUnicodeCharacter{2213}{\ensuremath{\mp}}
  \DeclareUnicodeCharacter{2217}{\ensuremath{\ast}}
  \DeclareUnicodeCharacter{221A}{\ensuremath{\sqrt{}}}
  \DeclareUnicodeCharacter{221E}{\ensuremath{\infty}}
  \DeclareUnicodeCharacter{2220}{\ensuremath{\angle}}
  \DeclareUnicodeCharacter{2227}{\ensuremath{\land}}
  \DeclareUnicodeCharacter{2228}{\ensuremath{\lor}}
  \DeclareUnicodeCharacter{2229}{\ensuremath{\cap}}
  \DeclareUnicodeCharacter{222A}{\ensuremath{\cup}}
  \DeclareUnicodeCharacter{222B}{\ensuremath{\int}}
  \DeclareUnicodeCharacter{2234}{\ensuremath{\therefore}}
  \DeclareUnicodeCharacter{2235}{\ensuremath{\because}}
  \DeclareUnicodeCharacter{2248}{\ensuremath{\approx}}
  \DeclareUnicodeCharacter{2260}{\ensuremath{\neq}}
  \DeclareUnicodeCharacter{2261}{\ensuremath{\equiv}}
  \DeclareUnicodeCharacter{2264}{\ensuremath{\leq}}
  \DeclareUnicodeCharacter{2265}{\ensuremath{\geq}}
  \DeclareUnicodeCharacter{226A}{\ensuremath{\ll}}
  \DeclareUnicodeCharacter{226B}{\ensuremath{\gg}}
  \DeclareUnicodeCharacter{2282}{\ensuremath{\subset}}
  \DeclareUnicodeCharacter{2283}{\ensuremath{\supset}}
  \DeclareUnicodeCharacter{2286}{\ensuremath{\subseteq}}
  \DeclareUnicodeCharacter{2287}{\ensuremath{\supseteq}}
  \DeclareUnicodeCharacter{22A5}{\ensuremath{\perp}}
  \DeclareUnicodeCharacter{22C5}{\ensuremath{\cdot}}
  \DeclareUnicodeCharacter{2295}{\ensuremath{\oplus}}
  \DeclareUnicodeCharacter{2297}{\ensuremath{\otimes}}
  % Greek letters
  \DeclareUnicodeCharacter{0391}{A}
  \DeclareUnicodeCharacter{0392}{B}
  \DeclareUnicodeCharacter{0393}{\ensuremath{\Gamma}}
  \DeclareUnicodeCharacter{0394}{\ensuremath{\Delta}}
  \DeclareUnicodeCharacter{0398}{\ensuremath{\Theta}}
  \DeclareUnicodeCharacter{039B}{\ensuremath{\Lambda}}
  \DeclareUnicodeCharacter{039E}{\ensuremath{\Xi}}
  \DeclareUnicodeCharacter{03A0}{\ensuremath{\Pi}}
  \DeclareUnicodeCharacter{03A3}{\ensuremath{\Sigma}}
  \DeclareUnicodeCharacter{03A6}{\ensuremath{\Phi}}
  \DeclareUnicodeCharacter{03A8}{\ensuremath{\Psi}}
  \DeclareUnicodeCharacter{03A9}{\ensuremath{\Omega}}
  \DeclareUnicodeCharacter{03B1}{\ensuremath{\alpha}}
  \DeclareUnicodeCharacter{03B2}{\ensuremath{\beta}}
  \DeclareUnicodeCharacter{03B3}{\ensuremath{\gamma}}
  \DeclareUnicodeCharacter{03B4}{\ensuremath{\delta}}
  \DeclareUnicodeCharacter{03B5}{\ensuremath{\varepsilon}}
  \DeclareUnicodeCharacter{03B6}{\ensuremath{\zeta}}
  \DeclareUnicodeCharacter{03B7}{\ensuremath{\eta}}
  \DeclareUnicodeCharacter{03B8}{\ensuremath{\theta}}
  \DeclareUnicodeCharacter{03B9}{\ensuremath{\iota}}
  \DeclareUnicodeCharacter{03BA}{\ensuremath{\kappa}}
  \DeclareUnicodeCharacter{03BB}{\ensuremath{\lambda}}
  \DeclareUnicodeCharacter{03BC}{\ensuremath{\mu}}
  \DeclareUnicodeCharacter{03BD}{\ensuremath{\nu}}
  \DeclareUnicodeCharacter{03BE}{\ensuremath{\xi}}
  \DeclareUnicodeCharacter{03C0}{\ensuremath{\pi}}
  \DeclareUnicodeCharacter{03C1}{\ensuremath{\rho}}
  \DeclareUnicodeCharacter{03C3}{\ensuremath{\sigma}}
  \DeclareUnicodeCharacter{03C4}{\ensuremath{\tau}}
  \DeclareUnicodeCharacter{03C5}{\ensuremath{\upsilon}}
  \DeclareUnicodeCharacter{03C6}{\ensuremath{\varphi}}
  \DeclareUnicodeCharacter{03C7}{\ensuremath{\chi}}
  \DeclareUnicodeCharacter{03C8}{\ensuremath{\psi}}
  \DeclareUnicodeCharacter{03C9}{\ensuremath{\omega}}
  % Dashes, spaces, quotes
  \DeclareUnicodeCharacter{2013}{--}
  \DeclareUnicodeCharacter{2014}{---}
  \DeclareUnicodeCharacter{2018}{`}
  \DeclareUnicodeCharacter{2019}{'}
  \DeclareUnicodeCharacter{201C}{``}
  \DeclareUnicodeCharacter{201D}{''}
  \DeclareUnicodeCharacter{2026}{\ldots}
  \DeclareUnicodeCharacter{00A0}{~}
  \DeclareUnicodeCharacter{2002}{\enspace}
  \DeclareUnicodeCharacter{2003}{\quad}
  % Miscellaneous symbols
  \DeclareUnicodeCharacter{00A9}{\textcopyright}
  \DeclareUnicodeCharacter{00AE}{\textregistered}
  \DeclareUnicodeCharacter{2122}{\texttrademark}
  \DeclareUnicodeCharacter{00B0}{\textdegree}
  \DeclareUnicodeCharacter{00B2}{\textsuperscript{2}}
  \DeclareUnicodeCharacter{00B3}{\textsuperscript{3}}
  \DeclareUnicodeCharacter{00B9}{\textsuperscript{1}}
  \DeclareUnicodeCharacter{00BC}{\textonequarter}
  \DeclareUnicodeCharacter{00BD}{\textonehalf}
  \DeclareUnicodeCharacter{00BE}{\textthreequarters}
  \DeclareUnicodeCharacter{2022}{\textbullet}
  \DeclareUnicodeCharacter{2032}{\ensuremath{\prime}}
  \DeclareUnicodeCharacter{2033}{\ensuremath{\prime\prime}}
  % Currency
  \DeclareUnicodeCharacter{20AC}{\texteuro}
  \DeclareUnicodeCharacter{00A3}{\textsterling}
  \DeclareUnicodeCharacter{00A5}{\textyen}
  % Check marks and crosses
  \DeclareUnicodeCharacter{2713}{\checkmark}
  \DeclareUnicodeCharacter{2714}{\checkmark}
  \DeclareUnicodeCharacter{2717}{\ensuremath{\times}}
  \DeclareUnicodeCharacter{2718}{\ensuremath{\times}}
  % Ballot boxes (task list checkboxes)
  \DeclareUnicodeCharacter{2610}{\mdcheckbox}
  \DeclareUnicodeCharacter{2611}{\mdcheckboxchecked}
  \DeclareUnicodeCharacter{2612}{\mdcheckboxchecked}
  % Box drawing / geometric
  \DeclareUnicodeCharacter{25A0}{\ensuremath{\blacksquare}}
  \DeclareUnicodeCharacter{25A1}{\ensuremath{\square}}
  \DeclareUnicodeCharacter{25CB}{\ensuremath{\circ}}
  \DeclareUnicodeCharacter{25CF}{\ensuremath{\bullet}}
  \DeclareUnicodeCharacter{25B2}{\ensuremath{\blacktriangle}}
  \DeclareUnicodeCharacter{25B6}{\ensuremath{\triangleright}}
  \DeclareUnicodeCharacter{25BC}{\ensuremath{\blacktriangledown}}
  \DeclareUnicodeCharacter{25C0}{\ensuremath{\triangleleft}}
  \DeclareUnicodeCharacter{2605}{\ensuremath{\bigstar}}
  % Superscript / subscript digits
  \DeclareUnicodeCharacter{2070}{\textsuperscript{0}}
  \DeclareUnicodeCharacter{2074}{\textsuperscript{4}}
  \DeclareUnicodeCharacter{2075}{\textsuperscript{5}}
  \DeclareUnicodeCharacter{2076}{\textsuperscript{6}}
  \DeclareUnicodeCharacter{2077}{\textsuperscript{7}}
  \DeclareUnicodeCharacter{2078}{\textsuperscript{8}}
  \DeclareUnicodeCharacter{2079}{\textsuperscript{9}}
  \DeclareUnicodeCharacter{2080}{\textsubscript{0}}
  \DeclareUnicodeCharacter{2081}{\textsubscript{1}}
  \DeclareUnicodeCharacter{2082}{\textsubscript{2}}
  \DeclareUnicodeCharacter{2083}{\textsubscript{3}}
  \DeclareUnicodeCharacter{2084}{\textsubscript{4}}
  \DeclareUnicodeCharacter{2085}{\textsubscript{5}}
  \DeclareUnicodeCharacter{2086}{\textsubscript{6}}
  \DeclareUnicodeCharacter{2087}{\textsubscript{7}}
  \DeclareUnicodeCharacter{2088}{\textsubscript{8}}
  \DeclareUnicodeCharacter{2089}{\textsubscript{9}}
\else
  \usepackage{newunicodechar}
  % Arrows
  \newunicodechar{←}{\ensuremath{\leftarrow}}
  \newunicodechar{↑}{\ensuremath{\uparrow}}
  \newunicodechar{→}{\ensuremath{\rightarrow}}
  \newunicodechar{↓}{\ensuremath{\downarrow}}
  \newunicodechar{↔}{\ensuremath{\leftrightarrow}}
  \newunicodechar{↕}{\ensuremath{\updownarrow}}
  \newunicodechar{⇐}{\ensuremath{\Leftarrow}}
  \newunicodechar{⇒}{\ensuremath{\Rightarrow}}
  \newunicodechar{⇔}{\ensuremath{\Leftrightarrow}}
  % Math
  \newunicodechar{≈}{\ensuremath{\approx}}
  \newunicodechar{≠}{\ensuremath{\neq}}
  \newunicodechar{≤}{\ensuremath{\leq}}
  \newunicodechar{≥}{\ensuremath{\geq}}
  \newunicodechar{±}{\ensuremath{\pm}}
  \newunicodechar{×}{\ensuremath{\times}}
  \newunicodechar{÷}{\ensuremath{\div}}
  \newunicodechar{∞}{\ensuremath{\infty}}
  \newunicodechar{∈}{\ensuremath{\in}}
  \newunicodechar{∉}{\ensuremath{\notin}}
  \newunicodechar{∑}{\ensuremath{\sum}}
  \newunicodechar{√}{\ensuremath{\sqrt{}}}
  \newunicodechar{∀}{\ensuremath{\forall}}
  \newunicodechar{∃}{\ensuremath{\exists}}
  \newunicodechar{∧}{\ensuremath{\land}}
  \newunicodechar{∨}{\ensuremath{\lor}}
  \newunicodechar{∩}{\ensuremath{\cap}}
  \newunicodechar{∪}{\ensuremath{\cup}}
  \newunicodechar{·}{\ensuremath{\cdot}}
  \newunicodechar{…}{\ldots}
  \newunicodechar{°}{\textdegree}
  \newunicodechar{✓}{\checkmark}
  \newunicodechar{✗}{\ensuremath{\times}}
  \newunicodechar{★}{\ensuremath{\bigstar}}
  \newunicodechar{☐}{\mdcheckbox}
  \newunicodechar{☑}{\mdcheckboxchecked}
  \newunicodechar{☒}{\mdcheckboxchecked}
\fi

% Task list checkbox commands
\newcommand{\mdcheckbox}{%
  \tcbox[on line,boxsep=0pt,left=0.5pt,right=0.5pt,top=0.5pt,bottom=0.5pt,%
    colback=white,colframe=black!70,boxrule=0.55pt,arc=0.7pt,%
    nobeforeafter,box align=base]{\phantom{\rule{6pt}{6pt}}}%
}
\newcommand{\mdcheckboxchecked}{%
  \tcbox[on line,boxsep=0pt,left=0.5pt,right=0.5pt,top=0.5pt,bottom=0.5pt,%
    colback=linktextblue,colframe=linktextblue,boxrule=0.55pt,arc=0.7pt,%
    nobeforeafter,box align=base]{%
    \makebox[6pt][c]{\rule{0pt}{6pt}\color{white}\fontsize{5.5pt}{5.5pt}\selectfont\checkmark}%
  }%
}

% Box drawing character approximations using LaTeX rules.
% Used by the Lua filter for box drawing chars in non-verbatim contexts.
\newcommand{\mdruleh}{\rule[0.35ex]{0.7em}{0.4pt}}
\newcommand{\mdrulev}{\rule[-0.2ex]{0.4pt}{1.1ex}}
\newcommand{\mdcornertl}{\rule[0.35ex]{0.35em}{0.4pt}\rule[-0.2ex]{0.4pt}{0.55ex}}
\newcommand{\mdcornertr}{\rule[-0.2ex]{0.4pt}{0.55ex}\rule[0.35ex]{0.35em}{0.4pt}}
\newcommand{\mdcornerbl}{\rule[-0.2ex]{0.4pt}{0.55ex}\rule[0.35ex]{0.35em}{0.4pt}}
\newcommand{\mdcornerbr}{\rule[0.35ex]{0.35em}{0.4pt}\rule[-0.2ex]{0.4pt}{0.55ex}}
\newcommand{\mdteel}{\rule[-0.2ex]{0.4pt}{1.1ex}\rule[0.35ex]{0.35em}{0.4pt}}
\newcommand{\mdteer}{\rule[0.35ex]{0.35em}{0.4pt}\rule[-0.2ex]{0.4pt}{1.1ex}}
\newcommand{\mdteet}{\rule[0.35ex]{0.7em}{0.4pt}\hspace{-0.35em}\rule[-0.2ex]{0.4pt}{0.55ex}}
\newcommand{\mdteeb}{\rule[-0.2ex]{0.4pt}{0.55ex}\hspace{-0.4pt}\rule[0.35ex]{0.7em}{0.4pt}}
\newcommand{\mdcross}{\rule[-0.2ex]{0.4pt}{1.1ex}\hspace{-0.4pt}\rule[0.35ex]{0.7em}{0.4pt}}

% Invisible formatting character declarations for pdfTeX safety net.
% These are already handled by the Lua filter, but declared here as fallback
% in case any slip through (e.g., from template or metadata).
\ifPDFTeX
  \DeclareUnicodeCharacter{200B}{}%  ZERO WIDTH SPACE
  \DeclareUnicodeCharacter{200C}{}%  ZERO WIDTH NON-JOINER
  \DeclareUnicodeCharacter{200D}{}%  ZERO WIDTH JOINER
  \DeclareUnicodeCharacter{200E}{}%  LEFT-TO-RIGHT MARK
  \DeclareUnicodeCharacter{200F}{}%  RIGHT-TO-LEFT MARK
  \DeclareUnicodeCharacter{200A}{}%  HAIR SPACE
  \DeclareUnicodeCharacter{202A}{}%  LEFT-TO-RIGHT EMBEDDING
  \DeclareUnicodeCharacter{202B}{}%  RIGHT-TO-LEFT EMBEDDING
  \DeclareUnicodeCharacter{202C}{}%  POP DIRECTIONAL FORMATTING
  \DeclareUnicodeCharacter{202D}{}%  LEFT-TO-RIGHT OVERRIDE
  \DeclareUnicodeCharacter{202E}{}%  RIGHT-TO-LEFT OVERRIDE
  \DeclareUnicodeCharacter{2060}{}%  WORD JOINER
  \DeclareUnicodeCharacter{FEFF}{}%  BOM / ZERO WIDTH NO-BREAK SPACE
  \DeclareUnicodeCharacter{2010}{-}%  HYPHEN
  \DeclareUnicodeCharacter{2011}{-}%  NON-BREAKING HYPHEN
  \DeclareUnicodeCharacter{2012}{--}% FIGURE DASH
  \DeclareUnicodeCharacter{2015}{---}% HORIZONTAL BAR
  \DeclareUnicodeCharacter{2027}{\ensuremath{\cdot}}% HYPHENATION POINT
  \DeclareUnicodeCharacter{2500}{\mdruleh}%  BOX HORIZ
  \DeclareUnicodeCharacter{2502}{\mdrulev}%  BOX VERT
  \DeclareUnicodeCharacter{250C}{\mdcornertl}%
  \DeclareUnicodeCharacter{2510}{\mdcornertr}%
  \DeclareUnicodeCharacter{2514}{\mdcornerbl}%
  \DeclareUnicodeCharacter{2518}{\mdcornerbr}%
  \DeclareUnicodeCharacter{251C}{\mdteel}%
  \DeclareUnicodeCharacter{2524}{\mdteer}%
  \DeclareUnicodeCharacter{252C}{\mdteet}%
  \DeclareUnicodeCharacter{2534}{\mdteeb}%
  \DeclareUnicodeCharacter{253C}{\mdcross}%
\fi


\RedeclareSectionCommand[
  beforeskip=1.25\baselineskip,
  afterskip=0.5\baselineskip
]{paragraph}

\RedeclareSectionCommand[
  beforeskip=1.0\baselineskip,
  afterskip=0.4\baselineskip
]{subparagraph}

\usepackage[most]{tcolorbox}

\definecolor{linktextblue}{HTML}{2E6DA4}
\definecolor{tableheaderbg}{HTML}{F0F2F4}
\definecolor{table-rule-color}{HTML}{999999}
\definecolor{calloutnotebg}{HTML}{EAF3FF}
\definecolor{calloutnoteborder}{HTML}{2B6CB0}
\definecolor{calloutsuccessbg}{HTML}{EAF6EC}
\definecolor{calloutsuccessborder}{HTML}{2F7D4A}
\definecolor{calloutwarningbg}{HTML}{FFF4DE}
\definecolor{calloutwarningborder}{HTML}{B87400}
\definecolor{callouttipbg}{HTML}{EAF7F0}
\definecolor{callouttipborder}{HTML}{2D7A46}
\definecolor{calloutcautionbg}{HTML}{FDEDE8}
\definecolor{calloutcautionborder}{HTML}{C2410C}
\definecolor{calloutimportantbg}{HTML}{FCEBEC}
\definecolor{calloutimportantborder}{HTML}{B42318}
\definecolor{calloutinfobg}{HTML}{F1EFFC}
\definecolor{calloutinfoborder}{HTML}{5D4E9D}
\definecolor{codeblockbg}{HTML}{F4F7FA}
\definecolor{codeblockborder}{HTML}{D9D9D9}
\definecolor{codelinenumber}{HTML}{7A7A7A}
\definecolor{pagenumbergray}{HTML}{777777}
\definecolor{inlinecodebg}{HTML}{F3F3F3}
\definecolor{inlinecodefg}{HTML}{000000}

\addtokomafont{pagenumber}{\color{pagenumbergray}}

\newtcbox{\inlinecodebox}{
  on line,
  box align=base,
  nobeforeafter,
  enhanced,
  colback=inlinecodebg,
  colframe=inlinecodebg,
  boxrule=0pt,
  arc=2pt,
  left=3pt,
  right=3pt,
  top=1pt,
  bottom=1pt
}

\newcommand{\inlinecode}[1]{%
  \inlinecodebox{%
    \textcolor{inlinecodefg}{%
      \normalfont\mdseries\ttfamily\fontsize{9.85pt}{11pt}\selectfont #1%
    }%
  }%
}

\DeclareRobustCommand{\mdhighlight}[1]{%
  \leavevmode
  \begingroup
  \setlength{\fboxsep}{1pt}%
  \colorbox{yellow!35}{\strut #1}%
  \endgroup
}

\providecommand{\st}[1]{%
  \leavevmode
  \begingroup
  \setbox0=\hbox{#1}%
  \dimen0=\wd0
  \rlap{\raise0.55ex\hbox to \dimen0{\leaders\hrule height 0.08ex\hfill}}%
  \box0%
  \endgroup
}

\newlength{\mdemojiheight}
\newcommand{\mdemojiscaledheight}{}
\begingroup\expandafter\expandafter\expandafter\endgroup
\expandafter\ifx\csname pdfsuppresswarningpagegroup\endcsname\relax
\else
  \pdfsuppresswarningpagegroup=1\relax
\fi
\DeclareRobustCommand{\mdemoji}[1]{%
  \begingroup
  \settoheight{\mdemojiheight}{X}%
  \edef\mdemojiscaledheight{\the\dimexpr 11\mdemojiheight/10\relax}%
  \raisebox{-.1ex}{\includegraphics[page=#1,height=\mdemojiscaledheight]{all-twemojis.pdf}}%
  \endgroup
}

\newsavebox{\pandocTableBox}

\newenvironment{tightblockquoteheading}{%
  \begin{mdframed}[
    rightline=false,
    bottomline=false,
    topline=false,
    linewidth=3pt,
    linecolor=blockquote-border,
    skipabove=0pt,
    skipbelow=0pt,
    innerleftmargin=7pt,
    innerrightmargin=0pt,
    innertopmargin=2.7pt,
    innerbottommargin=2.7pt
  ]%
  \begingroup
  \color{blockquote-text}
}{%
  \endgroup
  \end{mdframed}
}

\newenvironment{mdspoilerbody}{%
  \begin{list}{}{%
    \setlength{\leftmargin}{1.5em}%
    \setlength{\rightmargin}{0pt}%
    \setlength{\listparindent}{0pt}%
    \setlength{\itemindent}{0pt}%
    \setlength{\parsep}{0pt}%
    \setlength{\topsep}{0pt}%
    \setlength{\partopsep}{0pt}%
  }%
  \item\relax
}{%
  \end{list}
}

\AtBeginDocument{%
  \definecolor{shadecolor}{HTML}{F4F7FA}%
  \setlength{\FrameSep}{0pt}%
  \renewcommand{\theFancyVerbLine}{\sffamily\textcolor{codelinenumber}{\scriptsize\arabic{FancyVerbLine}}}%
  \fvset{numbersep=9pt,xleftmargin=2.3em}%
  \renewenvironment{Shaded}{%
    \begin{tcolorbox}[
      enhanced,
      breakable,
      colback=codeblockbg,
      colframe=codeblockborder,
      boxrule=0.45pt,
      arc=1mm,
      left=5pt,
      right=5pt,
      top=6pt,
      bottom=6pt,
      boxsep=0pt,
      before skip=8pt,
      after skip=8pt
    ]%
  }{\end{tcolorbox}}%
}


\newtcolorbox{successbox}{
  enhanced,
  breakable,
  colback=calloutsuccessbg,
  colframe=calloutsuccessborder,
  boxrule=0.7pt,
  arc=1mm,
  left=2mm,
  right=2mm,
  top=1.2mm,
  bottom=1.2mm,
  before skip=8pt,
  after skip=8pt
}

\newtcolorbox{notebox}{
  enhanced,
  breakable,
  colback=calloutnotebg,
  colframe=calloutnoteborder,
  boxrule=0.7pt,
  arc=1mm,
  left=2mm,
  right=2mm,
  top=1.2mm,
  bottom=1.2mm,
  before skip=8pt,
  after skip=8pt
}

\newtcolorbox{warningbox}{
  enhanced,
  breakable,
  colback=calloutwarningbg,
  colframe=calloutwarningborder,
  boxrule=0.7pt,
  arc=1mm,
  left=2mm,
  right=2mm,
  top=1.2mm,
  bottom=1.2mm,
  before skip=8pt,
  after skip=8pt
}

\newtcolorbox{cautionbox}{
  enhanced,
  breakable,
  colback=calloutcautionbg,
  colframe=calloutcautionborder,
  boxrule=0.7pt,
  arc=1mm,
  left=2mm,
  right=2mm,
  top=1.2mm,
  bottom=1.2mm,
  before skip=8pt,
  after skip=8pt
}

\newtcolorbox{importantbox}{
  enhanced,
  breakable,
  colback=calloutimportantbg,
  colframe=calloutimportantborder,
  boxrule=0.7pt,
  arc=1mm,
  left=2mm,
  right=2mm,
  top=1.2mm,
  bottom=1.2mm,
  before skip=8pt,
  after skip=8pt
}

\newtcolorbox{tipbox}{
  enhanced,
  breakable,
  colback=callouttipbg,
  colframe=callouttipborder,
  boxrule=0.7pt,
  arc=1mm,
  left=2mm,
  right=2mm,
  top=1.2mm,
  bottom=1.2mm,
  before skip=8pt,
  after skip=8pt
}

\newtcolorbox{infobox}{
  enhanced,
  breakable,
  colback=calloutinfobg,
  colframe=calloutinfoborder,
  boxrule=0.7pt,
  arc=1mm,
  left=2mm,
  right=2mm,
  top=1.2mm,
  bottom=1.2mm,
  before skip=8pt,
  after skip=8pt
}

\AtBeginDocument{%
  \hypersetup{%
    pdfcreator={},%
    pdfproducer={},%
    pdfinfo={%
      /CreationDate ()%
      /ModDate ()%
    }%
  }%
}
